% !TeX root = ../../Book5.tex
% 纲要标题：### 14.1 sl_2 的结构
% 纲要位置：工作记录/计划纲要.md，第 4171 行。

\section{\texorpdfstring{\(\mathfrak{sl}_2\)}{sl₂} 的结构}
\label{b5:sec:1401}


% 纲要范围：
%
% * 标准基与交换关系
% * Cartan 子代数的例子
% * 升降算子
%

在最小的非交换半单例子中，三个算子就能展示最高权、升降与完全可约的关系。本章固定复数域，直接从三个交换关系推导全部有限维表示；后面的根系理论正是反复把较大的代数限制到这样的三元组上。

\subsection{标准基与交换关系}
\label{b5:subsec:1401:01}

在 \(\mathfrak{sl}_2(\mathbb C)\) 中取
\[
e=\begin{pmatrix}0&1\\0&0\end{pmatrix},\qquad
f=\begin{pmatrix}0&0\\1&0\end{pmatrix},\qquad
h=\begin{pmatrix}1&0\\0&-1\end{pmatrix}.
\]
它们构成一组基，直接相乘得到
\[
[h,e]=2e,\qquad [h,f]=-2f,\qquad [e,f]=h.
\]
因此在一个向量空间上给出 \(\mathfrak{sl}_2\) 表示，等价于给出满足这三条关系的算子 \(E,F,H\)。为使公式容易阅读，以下仍将作用算子写为 \(e,f,h\)；乘积 \(ef\) 指先作用 \(f\) 再作用 \(e\)，括号仍为交换子。
\begin{proposition}\label{b5:prop:sl2-simple}
复李代数 \(\mathfrak{sl}_2(\mathbb C)\) 是单李代数。记矩阵单位为 \(E_{ij}\)，取基 \(e=E_{12}\)、\(f=E_{21}\)、\(h=E_{11}-E_{22}\)。其 Killing 型满足
\[
\kappa(h,h)=8,\qquad \kappa(e,f)=\kappa(f,e)=4,
\]
其余基向量配对为零。
\end{proposition}
\begin{proof}
任一理想在 \(\operatorname{ad}h\) 下不变；该算子的三个特征值为 \(2,0,-2\)，对应直线分别是 \(\mathbb Ce,\mathbb Ch,\mathbb Cf\)。对一个非零理想中的元素，用这三个特征值的插值多项式作投影，可在理想中取得某个非零基向量。若取得 \(e\) 或 \(f\)，与另一基向量括号就得到 \(h\)，再得到全部三者；若先取得 \(h\)，则与 \(e,f\) 括号直接得到三者。因此理想只能为全部代数。
计算 Killing 型时，\(\operatorname{ad}h\) 的平方迹为 \(2^2+(-2)^2=8\)。算子 \((\operatorname{ad}e)(\operatorname{ad}f)\) 在 \(e,h,f\) 上的对角系数为 \(2,2,0\)，迹为四；反向积有同样的迹。两个同号根向量或 \(h\) 与 \(e,f\) 的乘积没有非零对角系数，故其余配对为零。
\end{proof}

\subsection{Cartan 子代数的例子}
\label{b5:subsec:1401:02}

令 \(\mathfrak h=\mathbb Ch\)。它是由对角矩阵组成的交换子代数，并给出伴随作用下的分解
\[
\mathfrak{sl}_2=\mathbb Cf\oplus\mathfrak h\oplus\mathbb Ce.
\]
对线性泛函 \(\alpha\in\mathfrak h^*\)、\(\alpha(h)=2\)，三项分别具有特征值函数 \(-\alpha,0,\alpha\)。\chapref{b5:ch:16}将从这样的分解建立 Cartan 子代数与根的一般定义；现在先看清这个例子中实际成立的性质。
\ProseParagraphBreak
与 \(h\) 对易的矩阵在 \(\mathfrak{sl}_2\) 中恰为 \(\mathfrak h\)。若
\(x=ae+bf+ch\) 满足 \([x,\mathfrak h]\subseteq\mathfrak h\)，则
\([x,h]=-2ae+2bf\) 迫使 \(a=b=0\)。所以 \(\mathfrak h\) 的中心化子和正规化子都等于自身；它的伴随算子又能同时对角化。这些性质是它作为 Cartan 子代数的具体表现，而不只是矩阵恰好为对角形。
\ProseParagraphBreak
选基的可逆对角矩阵会缩放 \(e,f\)，但不会改变关系的本质。例如取 \(a\ne0\)，令 \(e'=ae\)、\(f'=a^{-1}f\)、\(h'=h\)，仍满足相同三式。根向量的归一化可以变动；后面比较 Killing 对偶基和 Casimir 算子时，会把相应常数写出。

\subsection{升降算子}
\label{b5:subsec:1401:03}

\begin{definition}\label{b5:def:sl2-weight-space}
设 \(V\) 为复李代数 \(\mathfrak{sl}_2(\mathbb C)\) 的左模，记 \(h=\operatorname{diag}(1,-1)\)。对 \(\lambda\in\mathbb C\)，子空间
\[
V_\lambda=\{\,v\in V\mid hv=\lambda v\,\}
\]
称为关于 \(h\) 的权空间；若它非零，称 \(\lambda\) 为一个权。有限维时，再记 \(V_{[\lambda]}\) 为 \(h\) 的广义特征空间。
\end{definition}
由交换关系，
\[
eV_\lambda\subseteq V_{\lambda+2},\qquad
fV_\lambda\subseteq V_{\lambda-2}.
\]
因此称 \(e,f\) 为升算子与降算子。同样的包含也对广义特征空间成立，因为
\[
\Bigl(h-(\lambda+2)I\Bigr)^Ne=e(h-\lambda I)^N,\qquad
\Bigl(h-(\lambda-2)I\Bigr)^Nf=f(h-\lambda I)^N.
\]
此时尚不能假定 \(h\) 已经对角化，但线性算子的广义特征空间分解总可以使用。
\begin{lemma}\label{b5:lem:sl2-highest-vector-exists}
设 \(V\) 为非零有限维复 \(\mathfrak{sl}_2(\mathbb C)\) 模。记矩阵单位为 \(E_{ij}\)，取基 \(e=E_{12}\)、\(f=E_{21}\)、\(h=E_{11}-E_{22}\)。则存在 \(v\in V\setminus\{0\}\) 和 \(\lambda\in\mathbb C\)，使
\[
hv=\lambda v,\qquad ev=0.
\]
而且在每个有限维模上，算子 \(e,f\) 都幂零。
\end{lemma}
\begin{proof}
在 \(h\) 的有限个特征值中，选实部最大的 \(\lambda\)。此时 \(\lambda+2\) 不是特征值，所以 \(eV_{[\lambda]}=0\)。在这个非零广义特征空间中取一个非零特征向量，即得所求 \(v\)。
对幂零性，\(e^N\) 将 \(V_{[\mu]}\) 映入 \(V_{[\mu+2N]}\)。特征值集合有限，取充分大的同一个 \(N\)，所有这些目标都为零。因此 \(e^N=0\)。对 \(f\) 用 \(\mu-2N\) 同理。
\end{proof}
这个论证同时解释了复数域与有限维性的作用：需要存在特征值，并需要权不能无限地向同一方向延伸。无限维最高权模中，降算子通常就不幂零。
